\documentclass[a4paper]{article}

\usepackage{graphicx}
\usepackage{amsmath,amssymb}
\usepackage{minted}
\usepackage{multirow}
\usepackage{float}
\usepackage{todonotes}
\usepackage{forest}
\usepackage{xstring}
\usepackage{xcolor}
\usepackage[most]{tcolorbox}
\usepackage{xparse}
\usepackage{natbib}

\usetikzlibrary{graphs,graphdrawing}
\usegdlibrary{layered}

% for external graphviz
\input{/home/donkarlo/Dropbox/repo/nd_language_ai_project/src/nd_language_ai/formal/computer/programming/tex/latex/code_clips/external_graphviz.tex}

% for tags
\input{/home/donkarlo/Dropbox/repo/nd_language_ai_project/src/nd_language_ai/formal/computer/programming/tex/latex/part/control_sequence/kind/semtag/semtag.tex}

% for links
\input{/home/donkarlo/Dropbox/repo/nd_language_ai_project/src/nd_language_ai/formal/computer/programming/tex/latex/code_clips/seehere.tex}

\title{Generative}
\date{}

\begin{document}
   \maketitle

   \section{Definition}
      Generative model in ML:
      A model of how observed data are generated.

      Generative models are a class of computational models frequently used for classification. In machine learning, it typically models the joint distribution of inputs and outputs, such as P(X,Y), or it models how inputs are distributed within each class, such as P(X∣Y) together with a class prior P(Y). Because it describes a full data-generating process, a generative model can be used to draw new samples that resemble the observed data, a process often referred to as synthetic data generation.

      \seeherefooter{https://en.wikipedia.org/wiki/Generative_model}
      %\bibliographystyle{plainnat}
      %\bibliography{/home/donkarlo/Dropbox/repo/nd_shared_research/bibliography/bibliography.bib}
\end{document}